% Generated by build/make_tables.py from diversification_comparison.csv. Do not edit.

\begin{table}[tbp]
\centering
\small
\caption{The disclosed diversification effect against the proxy book, FY2016--FY2025}
\label{tab:results}
\begin{tabular}{@{}lrrrr@{}}
\toprule
 & \multicolumn{3}{c}{Diversification effect (\% of summed)} & Total VaR \\
\cmidrule(lr){2-4}
Year & Goldman & Resized yearly & Fixed since 2019 & error (\%) \\
\midrule
2016 & 41.7 & 53.9 & 49.3 & -21.0 \\
2017 & 39.3 & 47.6 & 46.1 & -13.7 \\
2018 & 41.2 & 44.8 & 43.2 & -6.1 \\
2019 & 41.7 & 45.1 & 45.1 & -5.8 \\
2020 & 44.4 & 49.6 & 45.5 & -9.4 \\
2021 & 39.0 & 48.1 & 45.8 & -14.9 \\
2022 & 45.5 & 43.6 & 42.4 & +3.3 \\
2023 & 41.1 & 37.7 & 34.5 & +5.7 \\
2024 & 43.6 & 39.1 & 32.3 & +8.0 \\
2025 & 43.0 & 37.9 & 31.9 & +9.1 \\
\midrule
Mean & 42.0 & 44.7 & 41.6 & \\
Std.\ dev. & 2.08 & 5.35 & 6.33 & \\
\bottomrule
\end{tabular}
\notetext{The diversification effect as a share of summed category VaR, by Goldman's own definition. \emph{Resized yearly} is the book recalibrated to each year's disclosed split; \emph{fixed since 2019} is sized once on 2019 and left alone. The final column is the error in the fixed book's total VaR against the disclosed total, which changes sign in 2022. The dollar level of the book is a normalisation, so that column measures the consequence of the diversification gap rather than a sizing error: Table~\ref{tab:sizing} shows the sizing is exact.}
\end{table}
